\documentclass[xcolor=x11names,compress,aspectratio=169]{beamer}
%\documentclass[xcolor=x11names,compress,aspectratio=43]{beamer}

%% General document %%%%%%%%%%%%%%%%%%%%%%%%%%%%%%%%%%
\usepackage{graphicx}
\usepackage{tikz}
\usepackage{amsmath}
\usepackage{amssymb}
\usepackage[utf8]{inputenc}
\usepackage{textpos}
\usepackage{booktabs}
\usepackage{listings}
\usepackage{auto-pst-pdf}

\usetikzlibrary{decorations.fractals, calc}
%%%%%%%%%%%%%%%%%%%%%%%%%%%%%%%%%%%%%%%%%%%%%%%%%%%%%%

%% Beamer Layout %%%%%%%%%%%%%%%%%%%%%%%%%%%%%%%%%%
\useoutertheme[subsection=false,shadow]{miniframes}
\useinnertheme{default}
%\usefonttheme{serif}
\usepackage{palatino}
\setbeamerfont{title like}{shape=\scshape}
\setbeamerfont{frametitle}{shape=\scshape}
\beamertemplatenavigationsymbolsempty


% Uncomment this line, if you want frame numbers on your slides
% (I disabled them since I had the progress bar at the header)
%\setbeamertemplate{footline}[frame number]

\definecolor{fmiBlue}{RGB}{11,128,145} % FSU Faculty blue
\definecolor{evbc}{RGB}{50,167,132} % EVBC Green
\definecolor{bgGray}{RGB}{230,230,230} % some gray if needed


\setbeamercolor*{upper separation line head}{bg=Snow4} %Snow4 comes from the x11names option of xcolor
%\setbeamercolor*{lower separation line head}{bg=evbc}  % obsolete - there is now a cool progress bar 
\setbeamercolor*{normal text}{fg=black,bg=white} 
\setbeamercolor*{alerted text}{fg=red} 
\setbeamercolor*{example text}{fg=black} 
\setbeamercolor*{structure}{fg=black} 
 
\setbeamercolor*{palette tertiary}{fg=black,bg=black!10} 
\setbeamercolor*{palette quaternary}{fg=black,bg=black!10} 

%%%%%%%%%%%%%%%%%%%%%%%%%%%%%%%%%%%%%%%%%%%%
% IMPORTANT
% Kevin:
% Currently the FMI blue (the light uni blue) is here.
% If you want for example the evbc green, you'd have to
% change the colors here. I will put this in a nice function
% at some point.
\setbeamercolor{title}{fg=evbc}
\setbeamercolor{structure}{fg=evbc}
\setbeamercolor{frametitle}{fg=evbc}
\setbeamercolor{block title}{fg=evbc, bg=bgGray}
\setbeamercolor{block title example}{fg=evbc}
\setbeamercolor{enumerate item}{fg=evbc}
\setbeamercolor{itemize item}{fg=evbc}
\setbeamercolor{item projected}{bg=evbc}

\setbeamerfont{block title}{size=\small}
\setbeamerfont{block body}{size=\scriptsize}

% Use those if you want. Just some shortcuts for 
% columns. Personally, I use minipages nowadays.
\renewcommand{\(}{\begin{columns}}
\renewcommand{\)}{\end{columns}}
\newcommand{\<}[1]{\begin{column}{#1}}
\renewcommand{\>}{\end{column}}

% Kevin:
% This is used for the code examples in the document.
% If used, set the frame option to 'fragile' (!!!)
\lstset{
	language=sh, 
	numbers=left, 
	numberstyle=\tiny\color{gray}, 
	backgroundcolor=\color{bgGray}, 
	basicstyle=\scriptsize\ttfamily,
	commentstyle=\tiny\color{Green4},
	keywordstyle=\color{Blue3},
	escapeinside=@@
}

% Kevin
% Appendix page number fix. More like an ugly hack.
% I was never able to create a "nice" appendix command, however, this one here works fine.
\newcommand{\beginbackup}{
   \newcounter{framenumbervorappendix}
   \setcounter{framenumbervorappendix}{\value{framenumber}}
}
\newcommand{\backupend}{
   \addtocounter{framenumbervorappendix}{-\value{framenumber}}
   \addtocounter{framenumber}{\value{framenumbervorappendix}} 
}

%%%%%%%%%%%%%%%%%%%%%%%%%%%%%%%%%%%%%%%%%%%%%%%%%%
% Kevin: Change this, if you want.
\graphicspath{{figures/}}
%

% Kevin:
% I used this slides for subsections a long time ago, however, at some point I figured that it is nicer to have those slides
% at the beginning of each section. Use this extra slides to give a short summary of the previous talk.
\AtBeginSection[]{
  \begin{frame}
  \vfill
  \centering
  \begin{beamercolorbox}[sep=8pt,center,shadow=true,rounded=true]{title}
    \usebeamerfont{title}\insertsection\par%
  \end{beamercolorbox}
  \vfill
  \end{frame}
}


% Kevin:
% small macro, used to build the progress bar
\def\insertframeratio{%
    \pgfmathparse{\insertframenumber/\inserttotalframenumber}%
}
%%%%%%%%%%%%%%%%%%%%%%%%%%
% Kevin:
% This is the progress bar in the headline.
% Change the textcolor if needed, change the width to 2pt if wanted
\addtobeamertemplate{headline}{}{%
  \insertframeratio
  \edef\myvar{\pgfmathresult}
  %\textcolor{evbc}{\rule{\myvar \paperwidth}{1pt}}
  \textcolor{evbc}{\rule{\paperwidth}{1pt}}
}

%%%%%%%%%%%%%%%%%%%%%%%%%%%
% Kevin
% This adds the FSU Logo to the
% bottom left. Changing the logo is easy
% adding a second logo (e.g. the evbc) is tricky
\addtobeamertemplate{footline}{}{%
\begin{textblock*}{100mm}(1em,-3.5em)
\includegraphics[scale=0.4]{evbc_cmyk.pdf}
\end{textblock*} 
\begin{textblock*}{100mm}(73em,-1em) 
\insertframenumber/\inserttotalframenumber
\end{textblock*}
}

%% Kevin:
%% I used this in order to show two logos at the same time.
%% as you can see, it is quite tricky in terms of positioning 
%% and size, but a nice try-and-error session usually
%% works out.
%\addtobeamertemplate{footline}{}{%
%\begin{textblock*}{100mm}(0.5em,-5em)
%\includegraphics[scale=0.5]{figures/evbc_cmyk.pdf}
%\end{textblock*}
%\begin{textblock*}{100mm}(11em,-3em)
%\includegraphics[scale=0.025]{fsu_logo.jpg}
%\end{textblock*}
%}


% Kevin:
% Well, fill this out as needed
%\title{Theoretical and practical metagenomic approaches to viral discovery}
%\subtitle{}
%\author{Kevin Lamkiewicz, Manja Marz}
%\date{27.09.2018\\[1em]RNA Bioinformatics and High-Throughput Analysis\\[1em]Friedrich Schiller University Jena}
%\date{xx.xx.2018\\RNA Bioinformatics and High-Throughput Analysis\\[2em]\includegraphics[scale=0.05]{fsu_logo.jpg}}



\title{Theoretical and practical metagenomic approaches to viral discovery}
\subtitle{Practical Session: Dimension Reduction and Clustering}
\author{Kevin Lamkiewicz, Manja Marz}
\date{23.10.2019\\[1em]European Virus Bioinformatics Center}

\begin{document}

\begin{frame}
  \maketitle
\end{frame}

\begin{frame}[c]\frametitle{}
  \begin{figure}[htbp]
    \centering
    \includegraphics<1>[width=1\textwidth]{dim_redu_1.pdf}
    \includegraphics<2>[width=1\textwidth]{dim_redu_2.pdf}
  \end{figure}
\end{frame}


\section[PCA]{PCA with scikit-learn}
\begin{frame}[c, fragile]\frametitle{A quick example}
  \begin{lstlisting}[language=Python, showstringspaces=false]
from sklearn.decomposition import PCA
from sklearn.preprocessing import StandardScaler
import numpy as np 
import pandas as pd 
import matplotlib.pyplot as plt


data = [] 
target = [] 
target_names = list(class2id.keys())
with open('virus.csv', 'r') as inputStream: 
    ...@\pause@
df = pd.DataFrame(data= np.c_[data, [id2class[x] for x in target]], 
        columns= ['feature1','feature2','feature3','feature4'] + ['target'])
  \end{lstlisting}
\end{frame}

\begin{frame}[c, fragile]\frametitle{Normalize your data}
  \begin{lstlisting}[language=Python, showstringspaces=false]
features = ['feature1','feature2','feature3','feature4']

# Here we seperate features from the target column
x = df.loc[:, features].values
y = df.loc[:,['target']].values

# Standardizing the only the features
x = StandardScaler().fit_transform(x)
  \end{lstlisting}
\end{frame}


\begin{frame}[c, fragile]\frametitle{Applying PCA}
  \begin{lstlisting}[language=Python, showstringspaces=false]    
pca = PCA(n_components=2)
principalComponents = pca.fit_transform(x)
principalDf = pd.DataFrame(data = principalComponents
         , columns = ['principal component 1', 'principal component 2'])
finalDf = pd.concat([principalDf, df[['target']]], axis = 1)
  \end{lstlisting}
\end{frame}

\begin{frame}[c]\frametitle{Other Application of PCAs}
  \begin{block}{Dimension Reduction in Machine Learning}
    Normally, you would like to use PCA (or any other method for dimension reduction) to \textbf{speed up}
    your machine learning algorithm.
    Instead of learning many instances with over 700 features (e.g. our handwritten letter example from earlier),
    we can \textbf{reduce the number of features}, by only taking the most \textbf{important combinations} of features into account.
  \end{block}
\end{frame}

\section[Classification]{Virus Classification with PCA}

\begin{frame}[c]\frametitle{Exercise / Homework}
  \begin{block}{}
  Have a look at the \texttt{viral\_data.fasta}.\\
  There are some known viruses, but also some unknown viruses inside. Try to use k-mer frequencies as features 
  and cluster the sequences based on that.\\Can you roughly classify the unknown viruses?\\ \ \\
  {\tiny You do not have to do this in Python. If you are more familiar with, for example, R, feel free to use this as well.\\
  I am just a little bit Python addicted. :)}
  \end{block}
\end{frame}


% Backup Slides. Using this macro, you'd get a slide number for each
% backup slide without increasing the maximum slide numbers of the original presentation.
% However, for this, the framenumber has to be inserted - which isn't in the template by default
\beginbackup

\begin{frame}[c]\frametitle{Coffee Break}
  \begin{figure}[htbp]
    \centering
    \includegraphics[width=0.65\textwidth]{coffeebreak.png}
  \end{figure}
\end{frame}

\backupend

\end{document}